\documentclass[11pt,a4paper]{article}

% --- PAQUETES DE IDIOMA Y CODIFICACIÓN ---
\usepackage[utf8]{inputenc}
\usepackage[T1]{fontenc}
\usepackage[spanish,es-tabla,es-nodecimaldot,es-noshorthands]{babel}

% --- MATEMÁTICAS Y SÍMBOLOS ---
\usepackage{amsmath,amssymb,amsfonts}
\usepackage{mathtools}

% --- FORMATO Y DISEÑO DE PÁGINA ---
\usepackage[top=2.5cm, bottom=2.5cm, left=2.5cm, right=2.5cm]{geometry}
\usepackage{fancyhdr}
\usepackage{booktabs}
\usepackage{array}
\usepackage{enumitem}
\usepackage{titlesec}

% --- COLORES Y CÓDIGO ---
\usepackage[table]{xcolor}
\definecolor{pujblue}{RGB}{0, 51, 102}
\definecolor{darkblue}{RGB}{16, 44, 87}
\definecolor{codebg}{RGB}{248, 249, 250}
\definecolor{darkgray}{RGB}{60, 60, 60}

\usepackage{listings}
\usepackage[colorlinks=true, linkcolor=pujblue, citecolor=pujblue, urlcolor=pujblue]{hyperref}

% --- CONFIGURACIÓN DE LISTINGS (SOPORTE COMPLETO UTF-8 / ESPAÑOL) ---
\lstset{
    language=Python,
    backgroundcolor=\color{codebg},
    basicstyle=\ttfamily\footnotesize,
    keywordstyle=\color{pujblue}\bfseries,
    commentstyle=\color{gray}\itshape,
    stringstyle=\color{teal},
    breaklines=true,
    frame=single,
    rulecolor=\color{black!20},
    numbers=left,
    numberstyle=\tiny\color{gray},
    numbersep=8pt,
    tabsize=4,
    showstringspaces=false,
    literate=
      {á}{{\'a}}1 {é}{{\'e}}1 {í}{{\'i}}1 {ó}{{\'o}}1 {ú}{{\'u}}1
      {Á}{{\'A}}1 {É}{{\'E}}1 {Í}{{\'I}}1 {Ó}{{\'O}}1 {Ú}{{\'U}}1
      {ñ}{{\~n}}1 {Ñ}{{\~N}}1 {ü}{{\"u}}1 {Ü}{{\"U}}1
      {¿}{{?`}}1 {¡}{{!`}}1
}

% --- ENCABEZADOS Y PIES DE PÁGINA ---
\pagestyle{fancy}
\fancyhf{}
\fancyhead[L]{\small\color{darkgray}\textit{Pontificia Universidad Javeriana Cali --- Computación Científica}}
\fancyhead[R]{\small\color{darkgray}\textit{Taller 1: Factorización LU (Método de Crout)}}
\fancyfoot[C]{\thepage}
\renewcommand{\headrulewidth}{0.4pt}
\renewcommand{\footrulewidth}{0.4pt}

\titleformat{\section}{\Large\bfseries\color{pujblue}}{\thesection}{1em}{}[\titlerule]
\titleformat{\subsection}{\large\bfseries\color{darkblue}}{\thesubsection}{1em}{}
\titleformat{\subsubsection}{\normalsize\bfseries\color{darkblue}}{\thesubsubsection}{1em}{}

% ==============================================================================
% INICIO DEL DOCUMENTO
% ==============================================================================
\begin{document}

\begin{center}
    {\LARGE\bfseries\color{pujblue} Taller 1: Factorización LU (Método de Crout) y Aplicaciones}\\[0.3cm]
    {\large Computación Científica --- Pontificia Universidad Javeriana Cali}\\[0.1cm]
    {\normalsize Profesor: Juan Fernando Paz}
\end{center}
\vspace{0.3cm}

% ==============================================================================
% ALGORITMOS BASE IMPLEMENTADOS EN PYTHON
% ==============================================================================
\section{Implementación de los Algoritmos Base en Python}

Para resolver ambos ejercicios se implementaron desde cero el \textbf{Método de Crout} ($A = LU$, donde $L$ es triangular inferior y $U$ es triangular superior con $1$s en su diagonal principal), junto con los algoritmos de \textbf{Sustitución Hacia Adelante} ($Ly = b$) y \textbf{Sustitución Hacia Atrás} ($Ux = y$).

\begin{lstlisting}
import numpy as np
from scipy.linalg import lu_factor, lu_solve

np.set_printoptions(precision=4, suppress=True)

def factorizacion_crout(A):
    A = np.array(A, dtype=float)
    n = A.shape[0]
    
    L = np.zeros((n, n), dtype=float)
    U = np.eye(n, dtype=float)  # Diagonal con 1s
    
    for j in range(n):
        # 1. Construccion de la columna j de L
        for i in range(j, n):
            suma_L = 0.0
            for k in range(j):
                suma_L += L[i, k] * U[k, j]
            L[i, j] = A[i, j] - suma_L
        
        # 2. Construccion de la fila j de U
        if L[j, j] != 0.0:
            for k2 in range(j + 1, n):
                suma_U = 0.0
                for r in range(j):
                    suma_U += L[j, r] * U[r, k2]
                U[j, k2] = (A[j, k2] - suma_U) / L[j, j]
            
    return L, U

def sustitucion_adelante(L, b):
    L = np.array(L, dtype=float)
    b = np.array(b, dtype=float)
    n = len(b)
    y = np.zeros(n, dtype=float)
    
    for i in range(n):
        suma = 0.0
        for k in range(i):
            suma += L[i, k] * y[k]
        y[i] = (b[i] - suma) / L[i, i]
        
    return y

def sustitucion_atras(U, y):
    U = np.array(U, dtype=float)
    y = np.array(y, dtype=float)
    n = len(y)
    x = np.zeros(n, dtype=float)
    
    for i in range(n - 1, -1, -1):
        suma = 0.0
        for k in range(i + 1, n):
            suma += U[i, k] * x[k]
        x[i] = (y[i] - suma) / U[i, i]
        
    return x
\end{lstlisting}

% ==============================================================================
% EJERCICIO 1
% ==============================================================================
\section{Ejercicio 1: Cantidad de Productos Generados}

El sistema de ecuaciones que modela los requerimientos de acero ($b_1$), plástico ($b_2$) y aluminio ($b_3$) para producir $x_1, x_2, x_3$ unidades de los productos A, B y C es:
\begin{equation}
    \underbrace{\begin{pmatrix}
        2 & 1 & 1 \\
        1 & 3 & 1 \\
        1 & 1 & 2
    \end{pmatrix}}_{A}
    \underbrace{\begin{pmatrix}
        x_1 \\ x_2 \\ x_3
    \end{pmatrix}}_{x}
    =
    \underbrace{\begin{pmatrix}
        b_1 \\ b_2 \\ b_3
    \end{pmatrix}}_{b}
\end{equation}

\subsection{Actividad 1: Factorización LU de la Matriz de Coeficientes}
Al aplicar la función \texttt{factorizacion\_crout(A)} sobre la matriz de coeficientes $A$, se obtienen las matrices $L$ y $U$:
\begin{equation}
    L = \begin{pmatrix}
        2.0 & 0.0 & 0.0 \\
        1.0 & 2.5 & 0.0 \\
        1.0 & 0.5 & 1.4
    \end{pmatrix}, \qquad
    U = \begin{pmatrix}
        1.0 & 0.5 & 0.5 \\
        0.0 & 1.0 & 0.2 \\
        0.0 & 0.0 & 1.0
    \end{pmatrix}
\end{equation}
Verificando el producto matricial se comprueba que $L \cdot U = A$ de forma exacta.

\subsection{Actividad 2: Cantidad de Productos para los Tres Escenarios}
Para cada vector de disponibilidad $b$, se resuelve primero $Ly = b$ (sustitución hacia adelante) y luego $Ux = y$ (sustitución hacia atrás):

\begin{table}[h!]
    \centering
    \caption{Resultados de producción para los tres escenarios de materia prima.}
    \begin{tabular}{llcc}
        \toprule
        \textbf{Escenario} & \textbf{Vector de Recursos $b$} & \textbf{Vector Intermedio $y$} & \textbf{Producción $x = [x_1, x_2, x_3]^T$} \\
        \midrule
        Escenario 1 & $[100, 120, 90]^T$ & $[50.0000, 28.0000, 18.5714]^T$ & $[28.5714, 24.2857, 18.5714]^T$ \\
        Escenario 2 & $[80, 100, 110]^T$ & $[40.0000, 24.0000, 41.4286]^T$ & $[11.4286, 15.7143, 41.4286]^T$ \\
        Escenario 3 & $[150, 160, 140]^T$ & $[75.0000, 34.0000, 34.2857]^T$ & $[44.2857, 27.1429, 34.2857]^T$ \\
        \bottomrule
    \end{tabular}
\end{table}

\begin{itemize}[leftmargin=*]
    \item \textbf{Escenario 1:} Producto A $= 28.5714$, Producto B $= 24.2857$, Producto C $= 18.5714$ unidades.
    \item \textbf{Escenario 2:} Producto A $= 11.4286$, Producto B $= 15.7143$, Producto C $= 41.4286$ unidades.
    \item \textbf{Escenario 3:} Producto A $= 44.2857$, Producto B $= 27.1429$, Producto C $= 34.2857$ unidades.
\end{itemize}

\subsection{Actividad 3: Explicación de Eficiencia Computacional (LU vs. Eliminación Gaussiana)}
La factorización LU resulta mucho más eficiente cuando se deben resolver múltiples sistemas lineales $Ax = b^{(k)}$ que comparten la misma matriz de coeficientes $A$:
\begin{enumerate}[leftmargin=*]
    \item \textbf{Con Eliminación Gaussiana desde cero:} Triangularizar la matriz aumentada $[A \mid b]$ requiere un costo de $\mathcal{O}\left(\frac{2}{3}n^3\right)$ operaciones. Si se tienen $m$ escenarios independientes, habría que repetir exactamente las mismas transformaciones sobre la matriz $A$ un total de $m$ veces, con un costo acumulado de $m \times \mathcal{O}\left(\frac{2}{3}n^3\right)$.
    \item \textbf{Con Factorización LU:} La descomposición $A = LU$ se calcula \textbf{una sola vez} con costo $\mathcal{O}\left(\frac{2}{3}n^3\right)$. Posteriormente, para cada nuevo vector $b$, únicamente se resuelven los dos sistemas triangulares $Ly = b$ y $Ux = y$ mediante sustitución hacia adelante y hacia atrás, lo cual cuesta apenas $\mathcal{O}(2n^2)$ operaciones por escenario. El costo total se reduce a $\mathcal{O}\left(\frac{2}{3}n^3\right) + m \times \mathcal{O}(2n^2)$.
\end{enumerate}

\subsection{Actividad 4: Solución con SciPy (\texttt{lu\_factor} y \texttt{lu\_solve}) y Comparación}
\begin{lstlisting}
lu_compacta, piv = lu_factor(A)

x1_scipy = lu_solve((lu_compacta, piv), np.array([100.0, 120.0, 90.0]))
x2_scipy = lu_solve((lu_compacta, piv), np.array([80.0, 100.0, 110.0]))
x3_scipy = lu_solve((lu_compacta, piv), np.array([150.0, 160.0, 140.0]))
\end{lstlisting}
\textbf{Comparación:}
\begin{itemize}[leftmargin=*]
    \item \textbf{Exactitud:} Tanto el método de Crout implementado manualmente como \texttt{lu\_factor}/\texttt{lu\_solve} de SciPy producen exactamente los mismos vectores solución:
    \begin{align*}
        x^{(1)}_{\text{SciPy}} &= [28.5714, \, 24.2857, \, 18.5714]^T \\
        x^{(2)}_{\text{SciPy}} &= [11.4286, \, 15.7143, \, 41.4286]^T \\
        x^{(3)}_{\text{SciPy}} &= [44.2857, \, 27.1429, \, 34.2857]^T
    \end{align*}
    \item \textbf{Diferencia práctica:} Frente a resolver cada sistema de manera independiente, el esquema en dos funciones (\texttt{lu\_factor} una sola vez y \texttt{lu\_solve} para cada escenario) encapsula la reutilización de la factorización en rutinas precompiladas de alto rendimiento (LAPACK), reduciendo el código requerido y garantizando estabilidad numérica mediante pivoteo parcial automático.
\end{itemize}

% ==============================================================================
% EJERCICIO 2
% ==============================================================================
\section{Ejercicio 2: Estimación de Emisiones Contaminantes en una Red de Sensores}

Se modela una red de $50$ sensores y $50$ fuentes de emisión mediante el sistema lineal $Ax = b$, con $A \in \mathbb{R}^{50 \times 50}$ y $b \in \mathbb{R}^{50}$.

\subsection{Código de Generación, Acondicionamiento y Resolución}
\begin{lstlisting}
def modificar_a_diagonal_dominante_estricta(A):
    A_mod = np.copy(A)
    n = A_mod.shape[0]
    for i in range(n):
        suma_fuera_diagonal = 0.0
        for j in range(n):
            if j != i:
                suma_fuera_diagonal += abs(A_mod[i, j])
        A_mod[i, i] = suma_fuera_diagonal + 1.0
    return A_mod

# Actividades 1, 2 y 3: Generacion de datos (semilla 42)
np.random.seed(42)
n = 50
A_sintetica = np.random.uniform(0.0, 1.0, size=(n, n))
A_dominante = modificar_a_diagonal_dominante_estricta(A_sintetica)
b = np.random.uniform(0.0, 100.0, size=n)

# Actividad 4: Factorizacion LU con Crout
L, U = factorizacion_crout(A_dominante)

# Actividad 5: Sustitucion hacia adelante y hacia atras
y = sustitucion_adelante(L, b)
x = sustitucion_atras(U, y)

# Actividad 6: Vector residuo y norma euclidiana
r = (A_dominante @ x) - b
norma_residuo = np.linalg.norm(r, 2)
\end{lstlisting}

\subsection{Desarrollo de las Actividades 1 a 6}
\begin{enumerate}[leftmargin=*]
    \item \textbf{Actividad 1 (Matriz $A$ aleatoria):} Se generó $A \in \mathbb{R}^{50 \times 50}$ con elementos distribuidos uniformemente en $[0, 1]$. Muestra de la submatriz inicial $A_{0:3, 0:3}$:
    \[
        A_{0:3, 0:3} = \begin{pmatrix}
            0.3745 & 0.9507 & 0.7320 \\
            0.9696 & 0.7751 & 0.9395 \\
            0.0314 & 0.6364 & 0.3144
        \end{pmatrix}
    \]
    \item \textbf{Actividad 2 (Modificación de la diagonal principal):} Se actualizó cada entrada diagonal según:
    \begin{equation}
        a_{ii} = \sum_{\substack{j=1 \\ j \neq i}}^{50} |a_{ij}| + 1
    \end{equation}
    obteniendo la matriz con \textbf{dominancia diagonal estricta}:
    \[
        A^{\text{dom}}_{0:3, 0:3} = \begin{pmatrix}
            22.9217 & 0.9507 & 0.7320 \\
            0.9696 & 24.9467 & 0.9395 \\
            0.0314 & 0.6364 & 24.6006
        \end{pmatrix}
    \]
    Esta propiedad garantiza matemáticamente (por el Teorema de Gershgorin) que $\det(A) \neq 0$ y que todos los pivotes $l_{ii}$ del método de Crout sean distintos de cero sin requerir intercambios de filas.
    \item \textbf{Actividad 3 (Vector de mediciones $b$):} Se generó $b \in \mathbb{R}^{50}$ con valores en $[0, 100]$. Primeras 5 mediciones:
    \[
        b_{0:5} = [82.7519, \; 76.4528, \; 57.3529, \; 95.6047, \; 20.0475]^T
    \]
    \item \textbf{Actividad 4 (Factorización con Crout):} Se obtuvieron $L \in \mathbb{R}^{50 \times 50}$ y $U \in \mathbb{R}^{50 \times 50}$ ($u_{ii} = 1$), con un error de factorización de:
    \[
        \|A - LU\|_2 = 3.91 \times 10^{-15}
    \]
    \item \textbf{Actividad 5 (Resolución $Ly = b$ y $Ux = y$):} Las primeras 5 intensidades de emisión estimadas fueron:
    \[
        x_{0:5} = [2.4363, \; 1.9972, \; 1.2748, \; 2.6976, \; -0.3245]^T
    \]
    \item \textbf{Actividad 6 (Verificación mediante el residuo):} Calculando el vector residuo $r = Ax - b$, se obtuvo la norma euclidiana:
    \begin{equation}
        \|r\|_2 = \|Ax - b\|_2 = 6.69 \times 10^{-14}
    \end{equation}
    Como $\|r\|_2 \approx 0$ (del orden de la precisión de máquina $10^{-14}$), se comprueba que la solución $x$ hallada es numéricamente exacta.
\end{enumerate}

\subsection{Pregunta de Análisis}
\textbf{Pregunta:} \textit{Suponga que la red de sensores permanece igual, pero se realizan nuevas mediciones en otro momento (es decir, la matriz $A$ permanece constante pero cambia el vector $b$). Explique qué parte del procedimiento podría reutilizarse para resolver el nuevo sistema de manera más eficiente.}

\medskip
\noindent\textbf{Respuesta:}
\begin{itemize}[leftmargin=*]
    \item \textbf{Parte que se reutiliza:} Se reutilizan directamente las matrices triangulares \textbf{$L$ y $U$} obtenidas en la factorización $A = LU$, ya que la matriz $A$ solo depende de la ubicación física y factores de dispersión constantes entre las fuentes y los sensores.
    \item \textbf{Operaciones para cada nueva medición $b_{\text{nuevo}}$:} Únicamente se deben ejecutar la \textbf{sustitución hacia adelante} ($Ly = b_{\text{nuevo}}$) y la \textbf{sustitución hacia atrás} ($Ux = y$).
    \item \textbf{Ventaja de eficiencia:} Factorizar la matriz $A$ de $50 \times 50$ requiere aproximadamente $\frac{2}{3}(50)^3 \approx 83,333$ operaciones, mientras que resolver los sistemas triangulares $Ly = b$ y $Ux = y$ requiere apenas $2(50)^2 = 5,000$ operaciones por cada nuevo vector de mediciones (un ahorro superior al $94\%$ del costo computacional).
\end{itemize}

\end{document}
